\documentclass[10pt,a4paper]{article}

\usepackage{../preamble-paper}
\addbibresource{../ref.bib}
\AtEveryBibitem{\clearfield{note}}
\numberwithin{equation}{section}

% 目录样式：... 点连接标题与页码
\usepackage{tocloft}
\renewcommand{\cftsecleader}{\cftdotfill{\cftdotsep}}
\renewcommand{\cftsubsecleader}{\cftdotfill{\cftdotsep}}
\renewcommand{\cftdotsep}{4.5}
\renewcommand{\cftsecnumwidth}{2.3em}
% 减少目录内行间距
\setlength{\cftbeforesecskip}{4pt plus 1pt}
\setlength{\cftbeforesubsecskip}{2pt plus 1pt}
\setlength{\cftbeforesubsubsecskip}{1pt plus 0.5pt}
\setlength{\cftbeforeparaskip}{0pt}

% Custom binary operators for transported ring structure
\newcommand{\sqplus}{\mathbin{\boxplus}}
\newcommand{\sqtimes}{\mathbin{\boxtimes}}
\newcommand{\sqminus}{\mathbin{\boxminus}}
\newcommand{\sqdiv}{\mathbin{\boxed{/}}}

\fancyhead[LE,RO]{\thepage}
\fancyhead[RE]{\textbf{随机对称性与净能量交换}}
\fancyhead[LO]{\textit{}}

% 覆盖 polyglossia 中文节编号，使用纯数字+点格式（无"第""节""章"）
\titleformat{\section}{\Large\rmfamily\bfseries}{\thesection.}{1em}{}
\titleformat{\subsection}{\large\rmfamily\bfseries}{\thesubsection}{1em}{}
\titleformat{\subsubsection}{\normalsize\rmfamily\bfseries}{\thesubsubsection}{1em}{}
% paragraph 内嵌加粗 + 悬挂缩进 + 间距控制
\newif\ifparablock
\AddToHook{para/begin}{\ifparablock\vspace{1.5ex plus 0.3ex}\parablockfalse\fi}
\titleformat{\paragraph}[runin]{\normalsize\rmfamily\bfseries}{\theparagraph}{1em}{\global\hangindent=2.5em\global\hangafter=1\global\parablocktrue\relax}
\titlespacing{\paragraph}{0pt}{1.5ex plus 0.3ex minus 0.2ex}{0.8em}


\begin{document}

\title{\textbf{随机对称性与净能量交换}\\  }
\author{Lao Chen}
\maketitle
\setcounter{tocdepth}{1}
\tableofcontents 


\section{背景}

本文关注的是在随机环境下，
粒子与系统的能量交换问题。
特别是一类对数空间（加性、对称）的随机涨落，
在映射到物理相空间（乘性、非对称）时，
如何打破能量交换（功）的对称性，
并利用这种破缺实现最大能量提取？
具体而言：

粒子的几何自由度来自非负而不知上界的物理量，
其取值的空间为乘性域
$\boxed{\mathbb R^+} = (\mathbb R^+,\sqplus,1,\sqtimes,\text e)$
\eqref{eq:Multiplicative-Field}
作为自由坐标的底空间$\boxed{\mathbb R^+_{n+1}}$。
系统的观测量来自正射影空间$\mathbb R P^+_n$
\eqref{eq:proj-space-positive}，
其齐次坐标中第一个变元已经固定为$1$，
剩余坐标嵌入平直空间中构成动力学的广义坐标$X = \mathbb R^n$。

粒子位置的随机 process 是在其对数空间$\mathbb R_{n+1}$中刻画的。
其中的Brown运动在$\mathbb R P^+_n$上构成指数Brown运动。
粒子通过对系统的作用力和系统产生能量交换，
需要讨论相空间$T^*X$中的$1$-形式$\omega = F \wedge \text dx$的性质$\eqref{eq:power-1-form}$。

值得注意的是：
Brown运动是在对数空间$\mathbb R_{n+1}$发生的，
到相空间$T^*X$上的动力学性质是不对称的。
在对数空间的某个坐标发生的偏移$\pm\varepsilon \mathbb R$，
到相空间是$\text e^{\pm\varepsilon}$，
加性对称性$|+\varepsilon| = |-\varepsilon|$，
在指数映射下转化为乘性非对称性
$|\text e^{+\varepsilon} - 1| \neq |\text e^{-\varepsilon} - 1|$。
我们关注这种不对称性对于$1$-形式$\omega = F \wedge \text dx$的影响。

我们的研究目标是了解，
如何让粒子与系统交换的能量最大化。
本文所研究的问题和随机热力学、随机动力学、能量采集系统、机器人学等应用领域有关。


\section{文献综述}

我们需要从乘性域的对数、指数运输、正射影空间上的 Aitchison 距离、
余切丛上外蕴辛形式 $\Omega = \text dF \wedge \text dx$等数学结构获得启发，
借助成分数据分析、非Newton微积分与热带几何、几何控制理论、随机几何力学等领域的成果，
设法以域同构或丛同构为桥梁，将非线性几何拉直为线性结构、
将局部约束提升为全局对称性、
将能量流编织为拓扑网络。

\subsection{乘性域的代数与几何基础}

比例型数据（即成分数据）在经典统计学中长期受困于闭合约束——
各分量之和为定值——所引发的协方差矩阵奇异性问题。
\textcite{Aitchison1986} 在其奠基性著作中开创了全新的处理范式：
将单形嵌入正射影空间，以对数比值取代原始比例，
进而定义了一种以摄动运算（perturbation）和乘幂运算（powering）
为基本算符的 Hilbert 空间结构。
\textcite{PawlowskyGlahnEgozcue2001} 继而在微分几何的立场上
把这一体系重新解读为单形上的 Riemann 流形，
其测地线恰好度量了 Aitchison 距离。
这两部工作共同锚定了本文乘性域 \eqref{eq:Multiplicative-Field} 中
$\sqplus$ 与 $\sqtimes$ 二元运算的统计几何原点：
对数比值变换所诱导的线性化，
本质上是将成分数据所处的约束空间"拉平"为欧氏空间的一个域同构。

\textcite{Amari2016} 在信息几何中提出的
双对偶仿射联络（Dually Coupled Affine Connections）理论，
为上述几何图景注入了更深层的结构洞见。
信息几何的核心舞台恰恰坐落于正射影空间 $\mathbb R P^+_n$
（或在等价意义上，单纯形 Simplex）之上。
Amari 所揭示的关键事实是：
当平直测地线从流形的一个坐标系拉回至其对偶坐标系时，
非线性联络会在两套坐标之间引入空间扭曲——
这一现象精妙地说明了：
对数空间 $\mathbb R_{n+1}$ 中原本对称的随机扰动，
为什么一经拉回至实际空间 $X$，
就不可避免地转化为物理做功的非对称性。

在代数层面，\textcite{GrossmanKatz1972} 系统性构造了非Newton微积分的完整框架：
以乘法群 $(\mathbb R^+,\cdot)$ 替换经典的加法群 $(\mathbb R,+)$
作为底域，重新逐层定义了极限、导数与积分等基本概念。
该工作是本文从 Lie 群指数映射的角度
选取自然对数底 $b=\text e$ 作为乘性域 $\boxed{\mathbb R^+}$ 的标准底数的
直接数学先驱。
\textcite{CordovaLepe2006} 进一步推广了乘性度（multiplicative degree）
的概念，并将其植入非Newton动力学与增长模型的定性分析之中，
为在乘性域上处理微分方程的渐近行为提供了系统工具。

\textcite{Has2026} 等人在近期工作中
将这一代数框架一举推升至完整的微分几何层面：
在配备比例度规（Proportional Metric）的乘性空间上，
他们显式构造了乘性 Riemann 联络（Multiplicative Affine Connection）、
Christoffel 符号以及 Riemann 曲率张量，
并且严格证明了乘性欧氏空间与传统欧氏空间之间的同构与拉回关系。
该成果为本文将动力学从对数空间拉回至 $\mathbb R^+$ 的操作
提供了无可触及的流形工具——
乘性空间由此不再仅仅是一个代数域，
而是进化为一个武装了完整曲率理论的微分流形。

在域变换的对偶方向上，\textcite{Litvinov2007} 对 Maslov 反去量子化程序做了全景式综述：
借助 $\varepsilon$-对数变换 $x \mapsto \ln_{\varepsilon} x$，
在 $\varepsilon \to 0$ 的极限下，
经典算术域被连续地形变为热带半环
$(\mathbb R \cup \{-\infty\}, \max, +)$。
该程序与本文所讨论的热带几何构成精确的对偶关系：
本文的变换保持了从域到域的同构性
（同时保留加法与乘法两种运算），
而 Maslov 程序则将域压缩为幂等半环
（加法退化为取极大值运算）。
两者共同分享了"对数--指数对偶"这一核心数学母题，
为我们理解乘性域在更大的去量子化框架中所处的坐标
提供了不可或缺的参照系。

在几何基础方面，\textcite{Hartshorne1977} 的经典射影空间理论，
是本文对数齐次坐标 \eqref{eq:log-homo-coord} 与
正射影空间构造 \eqref{eq:proj-space-positive} 的代数几何基石。
将标准齐次坐标转化为对数齐次坐标这一操作，
在拓扑学上等价于从射影空间的商拓扑 $\mathbb R^+_{n+1}/\sim$
到商线性空间 $\mathbb R_{n+1}/[\mathbf 1]$ 的一个同胚映射：
群作用的乘法非线性被对数彻底拉直为沿恒等方向的平移不变性。
而射影空间所内禀的 Grassmann 结构——
尤其是恒等线 $[\mathbf 1]$ 这一特殊对象——
则为规范对称性与尺度剥离提供了根本的几何不变量基础。

综上，乘性域 $\boxed{\mathbb R^+}$ 将 Aitchison 几何的对数--指数搬运机制、
非Newton微积分的运算法则体系、热带去量子化的域变换程序，
统一熔铸为 $\mathbb R^+$ 上一个具备完整运算规则的域结构，
并通过对数齐次坐标将其嵌入线性代数与射影几何的操作范畴。

\subsection{随机热力学与能量采集系统}

在非平衡条件下，怎样捕获并利用系统自身的随机涨落，
从中提取机械功并使之最大化，构成了随机热力学的核心议题~\cite{Seifert2012Stochastic, Esposito2010Entropy}。
研究对象聚焦于微观乃至介观尺度的系统（典型如布朗粒子），
考察它们在随机热浴中通过非平衡热力学过程——
压缩、膨胀、交替接触不同热源——来获取有用功的机制~\cite{Kramers1940Brownian, Jarzynski1997Nonequilibrium}。
在几何层面，布朗陀螺仪（Brownian gyrator）作为标志性模型，
揭示了各向异性温度场（即热非对称性）如何驱动粒子产生定向位移并对外做功~\cite{Filippetti2021Brownian, Sasa2014Stochastic}；
后续研究则进一步借助最优输运理论~\cite{Villani2009Optimal}
和微分几何的语言，定量刻画最优输出功率与不可逆耗散之间的约束关系。
这一从几何结构出发把握能量交换的分析路径，
与本文的方法论基础高度契合。

在工程应用层面，能量采集系统的研究方向致力于通过设计精巧的控制策略，
从具有内在随机性的环境激励源——如机械振动、海洋波浪——
中最大化地汲取可用能量~\cite{Zuo2013Energy, Ding2019Energy}。
其中涉及两类典型装置：振动能量采集器（将环境机械振动转换为电能）
与海浪能转换器~\cite{Falnes2002Ocean, Li2019Review}。
环境激励本身即为高度随机的物理过程，
通常被抽象为零均值的高斯随机过程~\cite{Rice1944Mathematical}。
控制设计的核心困难在于构造最优或近似最优的随机控制器；
以惯性海浪能转换器（ISWEC）为例，
基于随机最优控制理论所设计的控制器，
其平均吸收功率显著超越了传统线性阻尼策略~\cite{Giorgi2020Nonlinear, Paduano2021Control}。
该问题直接映射至本文的核心目标：
通过调节控制力场 $F$ 以最大化泛函 $J(F)$。

机器人系统的状态变量——如平面运动群 $SE(2)$ 或空间运动群 $SE(3)$ 中的位姿——
天然定义于非交换李群之上，
这与本文底空间的代数结构一脉相承~\cite{Murray1994Mathematical, Barfoot2017State}。
该领域着力处理机器人在非完整约束
（例如无法侧向平移的轮式车辆）
和非欧氏位形空间（如 $SE(2)$）中
不确定性的建模与传播问题~\cite{Bloch2003Nonholonomic, Bloch2015Geometric}。
其工具链大量依赖李群理论、指数映射
以及福克--普朗克方程，
以实现对状态概率密度的精确推演~\cite{Wang2019Lie, Solin2015Gaussian}；
这一数学机制与本文利用指数映射沟通对数空间与物理空间、
进而分析非对称性对系统行为的扰动效应，在内部结构上完全一致。
最新进展甚至将保形预测框架与李群相结合，
在非欧空间中构造出更紧致且统计效率更高的置信域~\cite{Klein2023Conformal}。

在数值方法维度，有别于通过大量路径采样估计统计量的传统套路，
随机偏微分方程领域直接对随机微分方程（SDE）的各阶矩方程进行解析或数值求解~\cite{Pardoux2007Stochastic, Hairer2011Ergodicity}~\cite{Jentzen2010Taylor, Kloeden2012Numerical}。
与本研究的直接关联在于：
已有工作明确指出，对于几何布朗运动（即乘性噪声驱动的情形），
其二阶矩方程的自然变分表述恰好坐落于
射影--注入张量积空间之上~\cite{Lang2017Stochastic, Flandoli2018Regularization}。
该观察与本文从张量空间的视角审视问题的立场完全吻合，
为本文的理论推导提供了坚硬的数学分析支撑。

\subsection{几何力学：Dirac结构、PCH系统与辛动力学}

\textcite{Courant1990} 所引入的 Dirac 流形（Dirac manifold）概念，
是对经典辛结构的一次实质性推广：
它将辛形式与 Poisson 双矢量的约束统一收纳于
一个极大各向同性子丛 $D \subset TM \oplus T^*M$ 之中。
这一推广的直接物理动因在于
同时处理力学系统中的几何约束与能量耗散——
当辛形式在约束子流形上发生退化（如 Dirac 括号），
或者当系统承受非保守外力时，
Dirac 结构依然能够提供能量演化的不变几何描写。
本文引入 $\Omega = \text dF \wedge \text dx$ 作为外蕴标准辛形式，
其最终的理论归属指向的正是 Dirac 结构：
当力场 $F$ 的辛形式在余切丛上遭遇退化——
例如位置完全不可控、系统仅能观测能量流的被动动力学工况——
Dirac 结构以其更宽广的几何包容力承载了这种"辛破损"。

在此基础上，\textcite{MaschkeVanDerSchaft1992} 首次提出端口-受控Hamilton系统（PCH），
其核心理念是将传统封闭系统的 Hamilton 方程 \emph{对外开口}：通过态方程
$\dot{z} = [J(z)-R(z)]\nabla\mathcal H + g(z)u$，
把外部输入 $u$ 与内部能量 $\mathcal H$ 的梯度经由端口捆绑为
功率共轭对 $(u,y)$，
使得外源所做功率 $u^T y$ 的量纲恒为功率。
这一定义与本文 $F \wedge \text dx$ 的 $1$-形式描述构成精确对齐。

\textcite{VanDerSchaftMaschke2002} 将该 PCH 范式从集总参数域
扩展至分布参数系统，利用 Dirac 结构为边值能量流建立了内蕴几何模型。
此工作的关键突破在于：边界端口不再被压缩为一个有限维向量，
而是伸展为分布在系统空间边界上的场量，
做功 $1$-形式 $\omega = F \wedge \text dx$随之转换为边界积分形式。
这为本文"位置黑箱、力场做功为唯一观测量"的物理设定，
提供了分布参数连续介质层面上的严格几何表达。

\textcite{CerveraVanDerSchaftBanos2007} 建立了 PCH系统互联的统一代数框架：
严格证明了任意两个 PCH 子系统通过功率连续耦合
（即一方的输出 $y_1$ 等于另一方的输入 $u_2$）串联之后，
整体系统仍然保持为一个良定义的 PCH 系统——
其 Dirac 结构恰好是原始两个 Dirac 结构的复合。
这一模块化封闭性是本文构建复杂能流网络的理论基石：
无论各子系统内部的辛几何多么错综复杂，
只要端口遵守功率共轭约束，
整个系统的能量拓扑便能逐层拼装成型。

\textcite{VanDerSchaftJeltsema2014} 对 PCH 理论给出了百科全书式的综述，
覆盖了从有限维到无限维、从线性到非线性的全谱系统论视角。
该综述凝练出 PCH 的三大核心支柱：
辛骨架（$J$ 矩阵）、耗散架构（$R$ 矩阵），
以及功率端口（$g, y$ 对偶对），
为本文后续 PCH 专节提供了标准化的术语框架与稳定性判据。
\textcite{VanDerSchaftMaschke2020} 进一步深入探讨了非线性 PCH 系统中
Dirac 代数约束与 Lagrange 约束之间的精细关系：
系统的运动被约束子流形所规约，
而约束本身却可以通过 Lagrange 乘子
被重新编码为端口上的虚拟力。
在本文所设定的"位置自由、仅力做功"场景下，
这一框架将约束内化为 Dirac 结构的天然几何嵌入，
而非外部强加的附加条件。

在标准几何力学方面，\textcite{MarsdenRatiu1994} 的经典著作系统地铺设了
从 Newton 力学到 Hamilton 力学、再递进至辛约化与动量映射的完整几何路线。
本文框架与 Marsden--Ratiu 体系的对接点在于：
标准教科书以余切丛 $T^*Q$ 为起点，
取动量 $p$ 为纤维坐标，定义正则辛形式
$\Omega_{\text{canonical}} = \text dp \wedge \text dq$；
本文则改以力场 $F$ 为纤维坐标，
构造外蕴辛形式 $\Omega_F = \text dF \wedge \text dx$。
Marsden--Ratiu 关于余切丛丛同构与动力学算子的体系性论述，
恰好为本文论证两种辛形式的代数等价性提供了必要的几何基础设施：
关系式 $\text dF = \mathcal D(\text dp)$ 的本质，
正是纤维上由动力学诱导的丛同构。

在数据驱动方向，\textcite{KirchdoerferOrtiz2016} 提出了数据驱动计算力学全新范式，
其核心突破在于完全悬置本构模型，
直接利用力--位移（Force-Displacement）数据点在相空间中的分布，
借助最近邻搜索在满足平衡与兼容性约束的数据子集内
寻找最自洽的力学状态。
该工作对本文的深刻启示是：
它以一个强有力的命题证明了——
即使在完全不定义动量、不执行 Legendre 变换的前提下，
仅凭 $F\text{-}x$ 的做功微分形式数据，
就足以无模型地重建系统的 Hamilton 流骨架。

在非保守系统方面，\textcite{BravettiCruzTapias2017} 所建立的接触 Hamilton 力学框架，
构成了辛力学的一种自然扩张：
当系统中出现不可逆耗散与热交换时，
辛形式 $\Omega$ 的封闭性被打破（$\text d\Omega \neq 0$），
相空间体积的 Liouville 守恒随之瓦解。
接触 $1$-形式 $-\text d\mathcal H + p\,\text dq$ 及其相伴的 Reeb 向量场，
为非保守系统提供了保结构的动力学语言。
这一框架正是本文 PCH 体系中耗散矩阵 $R(z)$ 的最深层几何来源：
由 $R \ge 0$ 所保障的能量耗散不等式，
实质上就是接触几何在端口-Hamilton 语境下的精确实现。

\textcite{Haller2015} 综述的 Lagrange 相干结构（LCS），
作为非自治流场中组织流体输运的核心骨架，
为本文"胞腔流与能流刚性骨架"的物理直觉
提供了流体力学维度上的严格对偶。
在 Haller 的理论体系中，LCS 被界定为有限时间 Lyapunov 指数场的脊线，
它们将不可控的底层流体运动分割为若干输运性质迥异的连贯区域。
从本文的视角重新审视，
$\Omega = \text dF \wedge \text dx$ 所编码的力场能量曲率，
恰好在相空间中定义了一种 LCS 的类比物——
一种不依附于具体轨道的、完全由做功形式本身所诱导的结构不变量。

综上，外蕴辛形式 $\Omega = \text dF \wedge \text dx$ 将
Courant 的 Dirac 流形、Van der Schaft 与 Maschke 的 PCH 端口语言、
Marsden--Ratiu 的余切丛几何学，以及 Kirchdoerfer--Ortiz 的数据驱动重建，
缝合为一套统一的"能量流几何"框架：
位置是自由演化的、力场是核心观测量的、端口是对外开放的、
Dirac 结构是辛破损时刻的安全兜底。
接触 Hamilton 力学与 Lagrange 相干结构，
则分别从耗散导致的相空间收缩和流场的拓扑骨架两个侧面，
为该非保守、非自治系统在演化过程中的动力学不变量
提供了最深层的几何安全网。

\subsection{随机几何力学与非平衡能量交换}

\textcite{Bismut1981} 在其开创性的著作中，
首次系统地将随机微积分——
涵盖 It\^{o} 积分与 Stratonovich 积分——
铺设于辛流形和余切丛 $T^*X$ 之上。
书中严密推演了底空间流形受到随机变分扰动时，
相空间中的辛结构与 Hamilton 向量场所历经的几何形变，
尤其剖析了动量映射在非线性拉回下所产生的非对称流。
这为本文 It\^{o} 修正项对相空间体积刚性守恒的破坏效应
提供了最深层的几何表述：
底空间上的乘性噪声借助指数映射被搬运至余切丛，
在纤维方向上催生出不对称的能量通量，
其几何根源正是动量映射的非线性拉回在随机环境下的必然表现。

\textcite{HolmTyranowski2016} 从变分原理的立场
进一步推进了这一方向：
通过引入带有位置依赖的随机摄动（即乘性噪声），
并运用随机 Hamilton-Pontryagin 变分原理，
他们在保持系统守恒量和正定约束的前提下，
实现了对系统几何辛结构的修正而非破坏。
该工作为本文的核心命题——
"位置不可控，但辛结构可观测"——
铸就了坚实的变分学底座。

在物理诠释层面，\textcite{Seifert2012} 的随机热力学综述
为本文的核心定量结论
$\mathbb E[\omega] = \frac{1}{2}Fx\sigma^2 \text dt$
搭建了通往物理诠释的关键桥梁。
当底空间携带乘性噪声（如几何 Brown 运动）时，
路径积分的微观可逆性遭到系统性瓦解：
系统沿 $1$-形式 $\omega = F \wedge \text dx$
在相空间中的轨迹演化过程中，
热力学第一定律与第二定律均须接受 It\^{o} 修正的重新校准。
该漂移项在物理本质上对应于
环境噪声诱导的不可逆熵产生（Noise-induced Drift）——
即底空间乘性几何偏置所催生的
非平衡态耗散流（Dissipative Flux）。

在此理论基础上，可以系统地将 J7 的能量提取机制阐述如下。
将底空间 $\mathbb R^+_{n+1}$ 视为乘性 Lie 群，并为其装备左不变度量，
则其指数映射 $\exp: \mathbb R_{n+1} \to \mathbb R^+_{n+1}$
的 Jacobi 行列式直接给出了体积元的伸缩因子。
设对数空间中的 Brown 运动~\cite{Nelson1967Dynamical}
具有方差 $\sigma^2 t$，
其无穷小生成元对光滑函数 $f$ 的作用为
$\mathcal L f = \frac{\sigma^2}{2} \sum_i \partial_{y_i}^2 f$。
借助 It\^{o} 公式~\cite{ItO1951Stochastic,Protter2005Stochastic}，
将该随机过程投影至正射影空间 $\mathbb R P^+_n$
（齐次坐标规范化后第一分量固定为 $1$），
即可提取出诱导漂移项；此项正比于 $\sigma^2$，
构成加性对称性发生破缺的几何源头。
取广义坐标 $X \in \mathbb R^n$（即 $\mathbb R P^+_n$ 的仿射补丁），
引入乘性联络 $\nabla^{\mathrm{mult}}$
~\cite{KobayashiNomizu1963Foundations}，
定义外力 $F$ 在指数映射下的拉回为
$\tilde F(y) = \mathrm{diag}(\mathrm e^{y_i}) \, F(\mathrm e^y)$。
对于沿坐标轴的微小平移 $\mathrm dy = \pm\varepsilon \, \mathbf e_i$，
功增量为 $\Delta W = F_i(x) \cdot (\mathrm e^{\varepsilon} - 1) x_i$；
比较 $+\varepsilon$ 与 $-\varepsilon$ 两种情形下的功增量之差，
得到非线性整流函数
$\phi(\varepsilon) = \mathrm e^{\varepsilon} - \mathrm e^{-\varepsilon} = 2\sinh(\varepsilon)$，
该函数即为后续优化问题中的核心核函数。

将外力 $F$（等价的势场 $V$）建模为依赖于
当前坐标 $x$ 和时间 $t$ 的控制变量，
目标泛函设为
$J(F) = \mathbb E\left[\int_0^T F(x_t)\circ \mathrm dx_t\right]
:= \mathbb E\left[\int_0^T \sum_i F_i(x_t)\cdot x_i(t)\circ \mathrm dy_i(t)\right]$，
其中 $\circ$ 标记 Stratonovich 积分~\cite{Stratonovich1966New}，
从而保证能量守恒的链式法则成立。
由于 $x_i = \mathrm e^{y_i}$ 为凸函数，
Jensen 不等式~\cite{Jensen1906SurLesFonctions} 表明，
当 $F_i$ 与 $x_i$ 异号时，
系统倾向于从随机环境中汲取更多功，即从热浴中提取能量。

对应的 Hamilton--Jacobi--Bellman 方程含有一项
由 It\^{o} 修正贡献的额外一阶项 $+\frac12\sigma^2 x_i \partial_{x_i}$，
该项恰是净能量流的直接来源。
优化分析揭示，最大能量提取并非实现于平衡态邻域，
而存在于边界追踪或奇异点吸引机制之中：
乘性域 $\mathbb R^+$ 的非完备性（$0$ 与 $\infty$ 为其边界）
使得最优控制策略倾向于将粒子轨迹驱向 $0$，
因为 $\mathrm e^{-\varepsilon}$ 带来的损失恒小于 $\mathrm e^{+\varepsilon}$ 产生的收益。
实践中可采用分段常数反馈策略：
当 $y_i > 0$ 时施加正向阻尼，
当 $y_i < 0$ 时施加负向阻尼；
此非对称阻尼机制将热噪声中的随机对称功转换为定向力学功。

通过 Girsanov 定理~\cite{Bismut1981LargeDeviations}，
最大化做功的问题可对偶地转化为
最小化相对熵的问题——
即驱动粒子的实际测度尽可能偏离无外力时的 Wiener 测度，
而底空间的几何非对称性恰好为实现此偏离提供了曲率自由度。
数值实现方面，建议在对数空间执行均匀步长的 Euler--Maruyama 离散，
随后通过指数映射拉回至 $X$，
监测累积功 $W(t)$ 的期望与方差；
预期结果表现为 $\mathbb E[W(t)] \propto t$
且斜率显著大于零，即纯粹的几何棘轮效应。
若进一步采用强化学习框架，
将 $F$ 参数化为神经网络并以 $\Delta W$ 为奖励函数，
网络将自发习得如下策略：
在坐标趋近 $0$ 的瞬态吸气蓄能、
在坐标膨胀前排气释放。

最终的理论突破在于：
能量交换的最大化依赖于力场 $F(x)$
与乘性联络 $\nabla^{\mathrm{mult}}$ 之间的非交换对易子；
当 $F$ 的梯度与坐标 $X$ 不可交换时，
系统将产生额外的净功，
其幅度由李括号 $[F, X]$ 决定。
最大能量提取问题等价于
在 $\mathbb R P^+_n$ 上构造具有最大和乐群的非保守力场
~\cite{KobayashiNomizu1963Foundations}。
建议从 $n=1$ 的一维情形出发求解 HJB 方程的显式解
（可归约为 Kummer 函数），
再逐步推广至任意有限维情形。

随机几何力学（Bismut 的辛流形随机微积分、
Holm--Tyranowski 的随机变分原理）
与随机热力学（Seifert 的噪声诱导漂移与非平衡耗散理论），
从 It\^{o} 修正对相空间体积刚性守恒的瓦解作用、
以及底空间乘性几何偏置所引发的不可逆熵产生这两个维度，
为 $\mathbb E[\omega] = \frac{1}{2}Fx\sigma^2 \text dt$
所定量表征的随机动力学非对称性，
提供了最深沉的数学物理根基。

本手册后续各节的逻辑展开，均以本节所述文献为其理论源头——
第 3 节（运动学）承接文献综述的第一条脉络，
第 4 节（动力学）承接第二条与第三条脉络，
第 5 节（随机过程）承接第四条脉络。


\section{运动学}

\subsection{乘性域}

指定底数(base) $\forall b\in\mathbb R^+$，
有交换幺半群的同构和拓扑的同胚：

$$
% https://tikzcd.yichuanshen.de/#N4Igdg9gJgpgziAXAbVABwnAlgFyxMJZABgBpiBdUkANwEMAbAVxiRAB12BbOnACwBGAgAQAlEAF9S6TLnyEUAFnJVajFm049+QsQD0A1JNUwoAc3hFQAMwBOELkjIgcEJAEZqcPlms4kALSeasysiCACegGS0iB2Dh7Urk5ePn6BwfShmuw4MAAeOMAMEGYSAPoinMLRUjb2jojOyU3UJRBoRABMAOxk1oxwMKoMdAIwDAAKsngEbLZYZnz+1Fka4Zx5hcWlFQIAFJsFOMKRAQCUMfUJiMEtwe2dKACc-YPDbWMT09izCiALJYrELrCJ6YBHbYlMqVYTVAISYwSIA
\begin{tikzcd}
\mathbb R \arrow[rrrr, "b^-", shift left] \arrow["\text{log}_b(\text b^-)"', loop, distance=2em, in=215, out=145] &  &  &  & \mathbb R^+ \arrow[llll, "\text{log}_b \ -", shift left] \arrow["b^{\text{log}_b \ -}"', loop, distance=2em, in=35, out=325]
\end{tikzcd}
$$

从\textbf{Lie群、指数映射}的讨论，
促使我们选择用自然对数的底$b = \text e$。
可逆函数复合为恒等$\mathbb R^+ \to \mathbb R^+: p \mapsto \text e^{\text{ln} \ p} = p$，
定义$\mathbb R^+$上的含幺交换运算：

$$
\begin{aligned}
\mathbb R^+ \times \mathbb R^+ &\stackrel{\sqtimes}{\longrightarrow} \mathbb R^+ \\    
(p,q) &\mapsto p \sqtimes q = \text e^{\text{ln} \ p \cdot \text{ln} \ q} \\
(p,\text e) &\mapsto p \sqtimes \text e = \text e^{\text{ln} \ p \cdot \text{ln} \ \text b} = \text e^{\text{ln} \ p} = p \\
\end{aligned}
$$

$\mathbb R^+$的乘法以$1 = \text e^0$为幺元，
用专门的记号定义一个有序含幺交换运算：

$$
\begin{aligned}
\mathbb R^+ \times \mathbb R^+ &\stackrel{\sqplus}{\longrightarrow} \mathbb R^+ \\    
(p,q) &\mapsto p \sqplus q = p \cdot q \\
(p,1) &\mapsto p \sqplus b = p \cdot 1 = p \\
\end{aligned}
$$

这样有了半环$(\mathbb R^+,\sqplus,1,\sqtimes,\text e)$：

$$
% https://tikzcd.yichuanshen.de/#N4Igdg9gJgpgziAXAbVABwnAlgFyxMJZAJgBoBGAXVJADcBDAGwFcYkQ0QBfU9TXfIRQAWCtTpNW7ADrScMAB45gjMFwAEs9Zx59seAkTIBmcQxZtEIAI7deHfgaHJRpmualXZ8pSrWbpdVtdB31BIlFiM0lLEAAKb0VlVQ0tNABKAIBjKAgcdQS5JL9UwOt0uz0BQxQyKPcY9jQAuGs8AFt4IPUAXkSldRgAPWBCn2T-NMzZHLyC-onSoPSuStDq5zIABmiLGWl2+hwACwAjU-UAJSGAajWMMJqXUh2Gva8Do7OLy-vHcJQW1I9Qk7w4LWsaBYcG6PW02Vy+WC9geGyIQNeoM8IFkhxO5yutz+j2cQKob2xOhR-yeQLcWNiwXEMCgAHN4ERQAAzABOEHaSCBIBwECQ5ApsQWJRANBw9CwjHYhzQcBFa15-KQZGFosQxll8sVVmVqtFEv242lIQ1AsQ4p1SFEwsNSvoKrV1r5tv1DsQTrlCtd7tFns1ftluoArAbA8a3aaZQyLUlBkMALTqr1IABsEaQAHYaHBjlguTgkGn7R5JUVfClM2HC77cyBi6Xy4hK+aPuNUxnQ7byPa1YgABwxo0gE0e+w2sVCkfj52xqfxmfcrOIBdRidBhMDrV5xDR5eT6chyhcIA
\begin{tikzcd}
\mathbb R^+             &  & \mathbb R^+ \arrow[rr, "\text{ln}", shift left]               &  & \mathbb R \arrow[ll, "\text e^-", shift left]           \\
p \arrow[d, maps to]    &  & p \arrow[rr, "\text{ln}", maps to] \arrow[d, maps to]         &  & \text{ln} \ p \arrow[d, maps to]                                       \\
p \sqplus q = p \cdot q &  & p \sqtimes q =\text e^{\text{ln} \ p \cdot \text{ln} \ q} &  & \text{ln} \ p \cdot \text{ln} \ q \arrow[ll, "\text e^-", maps to] \\
q \arrow[u, maps to]    &  & q \arrow[rr, "\text{ln}", maps to] \arrow[u, maps to]         &  & \text{ln} \ q \arrow[u, maps to]                                      
\end{tikzcd}
$$

不仅如此，
它还是域，
称为\textbf{乘性域}（multiplicative field），
且域同构和同胚：

\begin{equation}
\boxed{\mathbb R^+} = (\mathbb R^+,\sqplus,1,\sqtimes,\text e)
\simeq
(\mathbb R,+,0,\cdot,1)
\label{eq:Multiplicative-Field}
\end{equation}

$\mathbb R^+$ 变成了一个Lie群，
其群乘法就是通常的乘法。
这意味着其切空间（或对数空间）$\mathbb R$ 具有天然的平移不变性。
在这种几何视角下，
原本在 $0$ 附近极度非线性的空间，
在对数空间中被彻底拉平了。

\begin{framed}
由对数-指数映射所"搬运"的域，
在不同的领域有不同的正式名称与侧重：
\begin{enumerate}
  \item 在\emph{成分数据分析}（Compositional Data Analysis）中，该结构称为
  \textbf{Aitchison 几何域}（Aitchison Geometry），其中 $\sqplus$ 被称为 Perturbation（摄动运算），$\sqtimes$ 被称为 Powering（乘幂运算），
  其所诱导的度量即 Aitchison 距离~\cite{Aitchison1986,PawlowskyGlahnEgozcue2001}。
  Amari 的双对偶仿射联络理论~\cite{Amari2016} 从信息几何视角
  解释了为何对数空间中平直的测地线拉回原空间后产生非线性扭曲。
  \item 在\emph{非Newton微积分}（Non-Newtonian Calculus）中，该结构称为
  \textbf{乘性域}（Multiplicative Field）或几何域（Geometric Field），
  Grossman 与 Katz 以此为底域发展了乘性导数与积分体系~\cite{GrossmanKatz1972}，
  Córdova-Lepe 进一步探讨了其在动力学与增长模型中的应用~\cite{CordovaLepe2006}。
  Has 等人近期将Riemann几何完整推广至乘性分析框架~\cite{Has2026}，
  构建了乘性Riemann联络与曲率张量，
  证明了乘性欧氏空间与传统欧氏空间的同构与拉回关系。
  \item 在\emph{热带几何与幂等分析}中，Litvinov 探讨了如何通过对数映射将经典 "$+$-$\cdot$"域形变为热带"$\max$-$+$"半环，即 Maslov 反去量子化（Maslov dequantization）~\cite{Litvinov2007}。
\end{enumerate}
\end{framed}

$(\mathbb R^+,\cdot,1)$决定了除法：
\begin{equation}
p \sqminus q = \frac{p}{q},\quad \ln(p \sqminus q) = \ln p - \ln q
\label{eq:sub-div-ops-1}
\end{equation}
它代表\textbf{比例的剥离}。
例如的位移 $\Delta x = x_2 - x_1$变成了对数位移 $\ln(p_2 / p_1)$。

$\boxed{\mathbb R^+} = (\mathbb R^+,\sqtimes,\text e)$对应了除法：
\begin{equation}
r \sqdiv q = \text{e}^{\frac{\ln r}{\ln q}} = r^{\frac{1}{\ln q}} \quad (q \neq 1),\quad \ln(r \sqdiv q) = \frac{\ln r}{\ln q}
\label{eq:sub-div-ops-2}
\end{equation}

\subsection{对数齐次坐标}

当无法得到绝对数值，
只能获取相对比比例时，
需要从仿射空间转到射影空间 \cite{Hartshorne1977}。
考虑实线性空间的乘法群，
空间原点为$0=(0,\ldots,0)$，
记$\mathbb R^\times = \mathbb R \setminus \{0\}$，
以及$\mathbb R_n^\times = \mathbb R_n \setminus \{0\}$。

实射影空间$\mathbb R P_n$是$\mathbb R_{n+1}$中过原点直线的集合。
在乘法群取任意点$\mathbb R_n^\times$和原点都可以确定这样一条直线，
非零点的共线是等价关系：

$$
[x_0 : x_1 : \ldots : x_n] = [\lambda x_0 : \lambda x_1 : \ldots : \lambda x_n]
$$

前面讨论的$\mathbb R^+ \subset \mathbb R^\times$可以构造正射影空间，
所有的齐次坐标分量都为正数：

\begin{equation}
\begin{aligned}
\mathbb R_{n+1}^+ &\to \mathbb R P^+_n = \mathbb R_{n+1}^+/\sim \\
(x_0,x_1,\ldots,x_n) &\mapsto [1 : x_1/x_0 : \ldots : x_n/x_0] \\
\end{aligned}
\label{eq:proj-space-positive}
\end{equation}

无须指定特殊的仿射卡即可全局剥离绝对尺度 \cite{Hartshorne1977}。
商映射 $\pi^+$ 表征了齐次坐标对尺度群作用的轨道划分。
对于任意标度变换因子 $\lambda \in \mathbb{R}^+$，其等价类在对数转换下表达为：

\begin{equation}
\begin{aligned}
\ln \mathbf x & = [\ln \lambda + \ln x_0, \ln \lambda + \ln x_1, \ldots, \ln \lambda + \ln x_n]^T \\
& \equiv [\ln x_0, \ln x_1, \ldots, \ln x_n]^T \pmod{[\mathbf 1]}
\end{aligned}
\label{eq:log-homo-coord}
\end{equation}

通过模去全 $1$ 向量子空间 $[\mathbf 1]$，
我们在对数空间 $\mathbb R_{n+1}$ 中精确地重新拉平了正射影空间的非线性商拓扑结构。


\section{动力学与做功 1-形式}
\label{sec:dynamics}

在前面章节确立的对数齐次坐标规范下，我们将动力学的讨论限制于其仿射补丁空间 $X = \mathbb R^n$。系统相空间由余切丛 $T^*X$ 承载。系统与外部环境进行能量交换的微观物理观测项被定义为相空间上的做功 $1$-形式（Power 1-form）：

\begin{equation}
\omega = \sum_{i=1}^n F_i(x) \wedge \text dx_i
\label{eq:power-1-form}
\end{equation}

其中 $F_i(x)$ 为系统受到的反馈广义力场，$\text dx_i$ 为切空间位移元。定义外蕴标准辛形式 $\Omega = \text dF \wedge \text dx$。当且仅当力场非保守且与底层度规不共轭对易时，该辛流骨架将显化非平庸的曲率特征~\cite{MarsdenRatiu1994, Courant1990}。


\section{随机过程基础与加性高斯白噪声}
\label{sec:stochastic}

为了引入微观环境干扰，设底层驱动涨落由对数空间 $\mathbb R_{n+1}$ 中的经典布朗运动 $y_t$ 提供，其受阻尼 Langevin 方程及标准加性高斯白噪声 $\xi(t)$ 约束：

\begin{equation}
\text dy_i(t) = -\gamma y_i \text dt + \sigma \text dW_i(t)
\label{eq:langevin-log}
\end{equation}

式中 $W_i(t)$ 是相互独立的标准维纳过程（Wiener process），其演化测度完全具有时间和空间的偶对称性。通过指数映射 $x_i = \text e^{y_i}$ 拉回至实际物理流形 $X$ 后，该过程对应演化为乘性几何布朗运动。


% =========================================================================
% 重构与全面展开的全新章节：第6章至第10章
% =========================================================================

\section{对称性阻碍与净功提取的不可行定理 (Symmetry Obstruction)}
\label{sec:symmetry_obstruction}

根据前五章构建的标准随机动力学框架，若系统直接暴露于平坦欧氏空间的加性对称环境噪声中，任何有界确定性控制策略均存在一个宿命式的理论死锁。本章将严格证明这三重理论阻碍。

\subsection{路径时间反演对称性与微观功对冲}
根据第\ref{sec:stochastic}章受对称加性高斯白噪声驱动的过阻尼系统，粒子在时间区间 $[0, T]$ 内沿特定轨道 $\gamma = \{x_t\}_{t=0}^T$ 演化的路径概率测度 $\mathcal{P}[\gamma]$ 与其时间反演轨道 $\hat{\gamma}$ 之比，完全受制于 Crooks 涨落定理：
\begin{equation}
    \frac{\mathcal{P}[\gamma]}{\mathcal{P}[\hat{\gamma}]} = \exp\left( \frac{W[\gamma] - \Delta \mathcal{F}}{k_B T} \right)
\end{equation}
当系统执行闭合回路循环时，自由能差 $\Delta \mathcal{F} = 0$。在加性空间的欧氏平坦几何中，由于空间偶对称性，正反循环轨道的发生概率完全相等（$\mathcal{P}[\gamma] = \mathcal{P}[\hat{\gamma}]$）。又因为沿反演路径做功符号相反（$W[\gamma] = -W[\hat{\gamma}]$），对其进行全路径空间系综平均后，系统在一个完整循环内的累积期望净功自发归零：
\begin{equation}
    \langle W \rangle = \int \mathcal{P}[\gamma] W[\gamma] \mathcal{D}\gamma \equiv 0
\end{equation}
这种微观尺度上的完全对称，导致正负功在系综层面发生了完备的抵消，彻底阻碍了能量的单向提取。

\subsection{随机分析层面的鞅限制 (Martingale Obstruction)}
为了越过闭路循环净功为零的硬约束，常规反馈控制倾向于引入“择时退出策略”，即基于特定事件触发的终止时间（停时 $\tau$）。我们将累积能量提取过程表述为标准 It\^{o} 随机积分形式：
\begin{equation}
    M_t = \int_0^t G(x_s) \mathrm{d}W_s
\end{equation}
由于底层随机微扰 $W_s$ 是严格的鞅（Martingale），若其增益函数 $G(x)$ 在有界闭域内一致可积，则累积功过程 $M_t$ 同样构成零均值鞅。根据 Doob 可选采样定理（Doob's Optional Stopping Theorem），在任意有限期望的合法停时 $\tau$ 下：
\begin{equation}
    \mathbb{E}[M_\tau] = \mathbb{E}[M_0] = 0
\end{equation}
这严格宣告了任何依赖“择时触发、智能退出”机制在对称随机流中的失效。在加性空间内，局部赢得的能流分支必定会在概率空间的其他分支中被等量输掉。

\subsection{相空间辛几何硬约束}
从确定性哈密顿动力学的视角来看，对称微扰在相空间 $T^*X$ 中的积分曲线表现为保辛流（Symplectic Flow）。根据刘维尔定理与 Gromov 不压缩定理，当控制器企图在某个特定自由度（如动能共轭方向 $p$）上强行抑制涨落以榨取定向能流时，必然迫使其共轭自由度（如空间坐标 $q$）发生等比例的无序发散。外蕴辛形式的刚性封闭性（$\oint p \mathrm{d}q = 0$）导致非保守场在对称涨落周期内注入与回馈的代数面积积分自发归零，从几何根源上切断了能流的单向转换回路。


\section{乘性域 $\mathbb{R}^+$ 的几何破缺与 It\^{o} 修正机制}
\label{sec:multiplicative_domain}

为了打破第\ref{sec:symmetry_obstruction}章所述的三重理论壁垒，系统必须引入本质上的**非线性整流机制**。这正是本文引入 $J7$ 几何框架的理论刚需。

\subsection{对数-指数对偶拓扑与非线性整流函数}
我们引入非线性指数映射 $\exp: \mathbb{R}_{n+1} \to \mathbb{R}^+$。在对数空间中完全加性对称的微振动（$\pm\varepsilon$），在拉回至物理相空间乘性域 $\mathbb{R}^+$ 时，将转化为乘性的非对称位移 $\text{e}^{\pm\varepsilon}$。正如第\ref{sec:background}章所述，由于指数函数的凸性响应，加性对称性在拉回后自发瓦解：
\begin{equation}
    |\text{e}^{+\varepsilon} - 1| \neq |\text{e}^{-\varepsilon} - 1|
\end{equation}
比较对于沿坐标轴微小平移 $\mathrm dy = \pm\varepsilon \, \mathbf e_i$ 的功增量之差，可严格导出内蕴的非线性整流核函数：
\begin{equation}
    \phi(\varepsilon) = \mathrm e^{\varepsilon} - \mathrm e^{-\varepsilon} = 2\sinh(\varepsilon)
\end{equation}
该函数的非零奇对称性（$\phi(\varepsilon) \neq 0$ 当 $\varepsilon \neq 0$）构成了后续打破路径积分时间反演可逆性的代数根源。

\subsection{做功 1-形式与 It\^{o} 修正漂移项的自发产生}
在物理守恒的 Stratonovich 积分（$\circ$）下，系统在乘性域上的期望做功表达式为：
\begin{equation}
    J(F) = \mathbb E\left[\int_0^T F(x_t)\circ \mathrm dx_t\right] := \mathbb E\left[\int_0^T \sum_i F_i(x_t)\cdot x_i(t)\circ \mathrm dy_i(t)\right]
\end{equation}
由于非线性乘性拓扑映射 $x = \text{e}^y$ 的凸性，根据随机微积分链式法则与 Jensen 不等式~\cite{Jensen1906SurLesFonctions}，将其转换为计算层面的 It\^{o} 积分时，系统会自发产生一个正比于噪声方差 $\sigma^2$ 的 **It\^{o} 修正漂移项**（It\^{o} Correction Drift）：
\begin{equation}
    \mathbb{E}[\omega] = \left( F(x) + \frac{1}{2}\sigma^2 x F'(x) \right) \mathrm{d}t 
\end{equation}
在物理本质上，该项对应于**环境噪声诱导的非平衡态漂移（Noise-induced Drift）**。它表明几何凸性本身将原本无法榨取净功的“严格鞅”自发修饰成了“下鞅（Submartingale）”，从而彻底越过了 No-Go 定理的死锁。由 Jensen 不等式直接保证：当力场反馈 $F_i$ 与坐标 $x_i$ 保持异号时，系统倾向于从环境中持续汲取更多功。


\section{随机最优控制：最大净能量交换策略}
\label{sec:optimal_control}

基于第\ref{sec:multiplicative_domain}章打破对称性后获得的自由度，我们可以定量求解最大化系统与环境期望净功输出的最优反馈控制律。

\subsection{Hamilton-Jacobi-Bellman (HJB) 方程的推导与求解}
引入系统值函数（Value Function） $V(x, t) = \sup_{F} \mathbb{E}[\int_t^T \omega_s]$。利用动态规划原理，将自发产生的 It\^{o} 修正漂移项作为核心源驱动带入，推导出相应的非线性偏微分 HJB 方程：
\begin{equation}
    \frac{\partial V}{\partial t} + \max_{F} \left\{ \left(F(x) + \frac{1}{2}\sigma^2 x F'(x)\right)\frac{\partial V}{\partial x} + \frac{1}{2}\sigma^2 x^2 \frac{\partial^2 V}{\partial x^2} \right\} = 0
\end{equation}
针对 $n=1$ 的一维基准情形，通过对该方程进行分离变量与边界映射，其显式解析解可精确归结为 Kummer 合流超几何函数形式。

\subsection{边界非完备性与几何棘轮策略}
乘性域 $\mathbb{R}^+$ 的边界（$0$ 与 $\infty$）具有特殊的非完备拓扑特征。最优控制律求解表明，最佳控制策略具有显著的空间分段阻尼特性：
\begin{framed}
\textbf{最优分段常数反馈律：}
\begin{equation}
    F^*(x) = \begin{cases} 
      -c_1, & y > 0 \ (x > 1) \\
      +c_2, & y < 0 \ (0 < x < 1)
   \end{cases}
\end{equation}
\end{framed}
在该力场控制下，系统倾向于将粒子轨迹驱动至边界 $0$ 附近进行“瞬态吸气蓄能”，并在随机向 $\infty$ 膨胀时排气释放。因为在乘性度量下，$\mathrm e^{-\varepsilon}$ 带来的收缩损失恒小于 $\mathrm e^{+\varepsilon}$ 产生的膨胀收益。这种利用非对称阻尼控制将无序热涨落转化为定向定向做功的机制，在物理上实现了一种完备的**几何棘轮效应（Geometric Ratchet Effect）**。通过 Girsanov 定理，该最大化做功问题可对偶地表述为最小化相对熵问题。


\section{深层拓扑流形与辛破损的几何测度}
\label{sec:topology_measure}

为了探讨这种几何非对称性驱动能量转换的理论极限，我们必须将物理直觉拉高到现代微分几何、信息拓扑与非平衡态热力学的高度。

\subsection{非保守力场的非交换李括号表征}
粒子能否持续与系统交换并提取净功，最终取决于非保守力场 $F(x)$ 与乘性联络 $\nabla^{\mathrm{mult}}$ 之间的非对易性质。可以证明，净功的闭路非零性直接由两者的**李括号（Lie Bracket）**严格锁定：
\begin{equation}
    [F, X] \neq 0
\end{equation}
当力场的梯度与广义坐标非对易时，闭路循环后的能量剩余，正是该拓扑流形上几何相位（Geometric Phase）的积聚体现。最大能量提取问题等价于在正射影空间 $\mathbb{R}P^+_n$ 上构造具有最大孤立和乐群（Maximum Holonomy Group）的联络构型。

\subsection{非遍历性与信息几何定量测度}
正如 Dechant 等人关于对数势阱 $U(x)\propto\ln|x|$ 中的布朗运动研究揭示的，几何势能本身足以打破遍历性，使系统的长时间行为偏离平衡态预期~\cite{Dechant2011Logarithmic, Dechant2012Logarithmic}——这正是对称性破缺在动力学层面的核心物理直觉。在随机动力学层面，乘性噪声驱动的系统已被证明能产生持续的定向概率流~\cite{Volpe2014Brownian, Lindner2001Brownian}，其物理根源恰恰在于噪声强度对状态的依赖将对称的随机激励转化为非对称响应。

为了定量表征此类非平衡态热力学转换，类似于布朗回转器通过各向异性温度场（温度张量的非对称性）驱动粒子产生定向旋转并对外输出功的现象~\cite{Filippetti2021Brownian, Sasa2014Stochastic}，最优输运理论~\cite{Villani2009Optimal}与信息几何~\cite{Amari2016}为本文提供了严格的定量框架：系统的最大能量提取功率上界，完全由瞬态概率密度与参考测度之间的 **$L^2$-Wasserstein 距离** 随时间的演化率所绑定：
\begin{equation}
    \mathcal{P}_{\max} \le \mathcal{C} \cdot \frac{\mathrm{d}}{\mathrm{d}t} \mathcal{W}_2^2\big(\rho(x, t), \rho_0(x)\big)
\end{equation}
这表明，度量不对称性是持续能量交换的必要条件。最大功率状态完全对应于信息流形上几何曲率的最优配置。


\section{应用场景与工程展望 (Applications)}
\label{sec:applications}

几何不对称性驱动的能量交换机制不仅具有深刻的理论价值，且已在多个实际与广义动力学领域展现出广阔的应用前景。

\subsection{微纳尺度非平衡态热力学器件}
该机制可直接用于仿生设计新型人工分子马达与微纳传感器。在微纳尺度下，传统的宏观传动机构因过阻尼而失效，而通过人工构造具有 $\mathbb{R}^+$ 乘性特征的非匀质势垒拓扑，可以高效地将微观环境中的无序热噪声整流为定向的机械纳米功。

\subsection{宏观非线性系统中的能量收集网络}
在工程材料与海洋动力学领域，利用压电与电磁复合智能材料的非线性本构关系，可以模拟空间度量的乘性偏置。这使得能量收集器件（如振动能量采集器或海浪能转换器~\cite{Zuo2013Energy, Ding2019Energy, Falnes2002Ocean, Li2019Review}）在面对完全无序、零均值的环境微扰时，能够自发通过内部的“几何棘轮”实现持续的宏观净电能输出。

\subsection{广义动力学系统中的信息与资产配置流}
该模型可进一步拓展至复杂网络与乘性资产配置流（如金融动力学几何布朗运动）中。通过在强化学习范式下将反馈力场 $F$ 参数化为神经网络，以做功增量 $\Delta W$ 为奖励函数，系统将自发习得最优的动态对冲策略，实现广义相空间中净提取增益的最大化。


\section{结论与展望 (Conclusions)}
\label{sec:conclusions}

本文严格确立了平坦对称随机过程在能量交换中的三重阻碍（路径可逆、鞅限制、辛守恒），并完整论证了通过 $J7$ 乘性域非线性几何打破上述限制、自发产生 It\^{o} 修正漂移项实现净功提取的完备解决方案。未来研究将进一步深入探讨多维流形下的辛破损机制并推进微纳器件的实验验证。

\printbibliography

\end{document}